% =====================================================================
% main_points.tex -- FROZEN interpretation points for manuscript integration
% Generated 2026-08-12. Training program CLOSED; no further model runs.
%
% Every number here traces to a committed artifact in this directory or to
% results/ under the repo root. Do not restate a number from memory: cite the
% file. Superseded values are listed in Section "Retired claims" and must not
% reappear anywhere in the manuscript, appendix or rebuttal.
%
% PROVENANCE CLOSED 2026-08-12. All five previously prose-only statistics are now
% regenerated from committed results by scripts/reproduce_*.py through the shared
% scripts/statistics_utils.py, and indexed in results/final_provenance.json.
% [reproduced] = regenerated from committed result files; the reproduced value is
% authoritative where it differs from the historical one by bootstrap noise.
% =====================================================================

\section*{Evidence hierarchy (preserve this order)}

\begin{enumerate}
  \item \textbf{Primary mechanism.} Tree-guided encoding, $+8.32$ pp causal main
        effect on macro-F1 (95\% CI $[+3.56, +13.00]$), from the complete crossed
        $2\times2$. Objective main effect $-0.33$ (null); interaction $-0.11$ (null).
  \item \textbf{Independent corroboration.} The \texttt{raw\_full} sensitivity arm,
        which changes \emph{only} \texttt{encoder\_mode} and leaves the full objective
        intact, costs $-8.26$ pp with a 95\% paired interval of $[-17.93, -1.21]$ that
        excludes zero. This is \emph{independent corroboration} of the $2\times2$
        encoder effect from a different design on different datasets. \textbf{Do not
        pool the two estimates as though they were one experiment.}
  \item \textbf{Objective role.} Little aggregate macro-F1 contribution, but
        substantial false-alarm control (see below).
  \item \textbf{Robustness.} A broad hyperparameter plateau; nominal improvements are
        statistically indistinguishable from the default.
  \item \textbf{Protocol validity.} TON-IoT and CTU-13 both show strong
        split-dependent generalisation effects \emph{without} establishing a
        universally superior architecture.
  \item \textbf{Deployment.} Compact storage, strong MCC/FAR, lower analytical SOP
        count, but a slower current CPU implementation.
  \item \textbf{Early exit.} Conditional SOP savings, automatically disabled when
        validation cannot preserve quality, and \textbf{not} a latency improvement.
  \item \textbf{Quantisation.} Dataset-sensitive for the Raw MLP; no detectable
        Tree-to-Spike loss on the four-dataset FP32 sensitivity subset.
\end{enumerate}

% ---------------------------------------------------------------------
% PRIORITY BLOCK -- the six most important quantitative results.
% Provenance tags: [verified] recomputed from raw result files on 2026-08-12;
% [response record] taken from the reviewer-response record and not
% independently recomputable from committed artifacts.
% ---------------------------------------------------------------------
\section*{1. The predeclared one-point dense-twin non-inferiority objective FAILED}

\textbf{This is the manuscript's original central objective and it was not met. State it
plainly.} On the four IoT/IIoT/IoMT target-domain datasets the paired
Tree-to-Spike\,$-$\,DNN-Tree macro-F1 difference is $\mathbf{-0.55}$ pp, with
hierarchical-bootstrap 95\% CI $\mathbf{[-1.93, +1.03]}$
\textup{[reproduced: \texttt{scripts/reproduce\_noninferiority.py}; historical
$[-1.92,+1.03]$, the difference is bootstrap noise]}. The script emits
\texttt{noninferiority\_established: false}.

\begin{quote}
Because the lower bound of the 95\% interval extends below $-1$ pp, the predeclared
one-point non-inferiority criterion \textbf{was not established}.
\end{quote}

\noindent Any 2-pp margin analysis is \textbf{post hoc only} and must be labelled as such.
Per-dataset paired differences: CICIoT2023 $-0.22$, CICIoMT2024 $+0.18$, Edge-IIoTset
$-1.09$, TON-IoT $-1.08$.

\paragraph{$\!$Do not confuse these two numbers.} The formal target-domain result
($-0.55$ pp, NI not established) is distinct from the eleven-dataset \emph{descriptive}
comparison of $79.80$ vs $79.86$, a difference of only $\mathbf{-0.06}$ pp. Reporting the
$-0.06$ pp figure in place of the non-inferiority verdict would misrepresent the outcome.

% ---------------------------------------------------------------------
\section*{2. Eleven-dataset headline comparison and its statistical companion}

Means over 11 datasets $\times$ 5 seeds \textup{[verified]}:

\begin{center}\small
\begin{tabular}{lrrr}\toprule
model & macro-F1 & MCC & FAR (\%) \\\midrule
Vanilla SNN            & 76.46 & 0.836 & 4.11 \\
\textbf{Tree-to-Spike} & \textbf{79.80} & \textbf{0.905} & \textbf{3.15} \\
DNN-Tree               & 79.86 & 0.888 & 3.31 \\
TASNN adaptation       & 73.77 & 0.837 & 4.72 \\
\bottomrule\end{tabular}\end{center}

\noindent Tree-to-Spike gains $\mathbf{+3.33}$ pp over Vanilla SNN and $\mathbf{+6.03}$ pp
over the evaluated TASNN adaptation, while trailing DNN-Tree by $\mathbf{0.06}$ pp
descriptively. It holds the \emph{best} MCC and the \emph{lowest} FAR of the four.

\paragraph{Friedman and Nemenyi \textup{[verified]}.} Friedman
$p = 9.31\times10^{-4}$ (macro-F1), $9.80\times10^{-4}$ (accuracy),
$7.20\times10^{-3}$ (MCC). Nemenyi critical difference $\mathbf{CD = 1.41}$ at
$\alpha=0.05$ for $k=4$, $N=11$. Macro-F1 mean ranks: DNN-Tree $\mathbf{1.73}$,
Tree-to-Spike $\mathbf{2.18}$, Vanilla SNN $\mathbf{2.27}$, TASNN $\mathbf{3.82}$.

\begin{quote}
No rank separation is detected between Tree-to-Spike and either DNN-Tree or the Vanilla
SNN; the separation from the TASNN adaptation is detected.
\end{quote}

\noindent\textbf{Non-significance is not equivalence.} Do not report the absent
Tree-to-Spike/DNN-Tree separation as evidence of parity --- the formal
non-inferiority test in Section 1 is the instrument for that question, and it failed.

% ---------------------------------------------------------------------
\section*{3. The crossed $2\times2$: cell means, MCC effect, and additivity}

Cell means (macro-F1) \textup{[reproduced: \texttt{scripts/reproduce\_crossed\_2x2.py}; all four cells match the historical values exactly]}:

\begin{center}\small
\begin{tabular}{lrr}\toprule
 & CE objective & full objective \\\midrule
Raw TTFS encoder  & 77.86 & 77.59 \\
Tree TTFS encoder & \textbf{86.23} & \textbf{85.85} \\
\bottomrule\end{tabular}\end{center}

\noindent Effects with 95\% CIs \textup{[reproduced]}. Cell mapping: Raw+CE $=$ ladder \texttt{raw}, Raw+full $=$ \texttt{raw\_full}, Tree+CE $=$ \texttt{tree\_ce}, Tree+full $=$ shipped default:
\begin{itemize}
  \item Encoder main effect: $\mathbf{+8.32}$ macro-F1, CI $[+3.57, +13.08]$; and
        $\mathbf{+17.95}$ MCC, CI $[+2.18, +33.80]$. \textbf{Freeze the MCC effect too ---
        it is larger in relative terms than the macro-F1 effect and is currently the
        under-reported half of the primary result.}
  \item Objective main effect: $-0.33$, CI $[-1.46, +0.66]$ (null).
  \item Interaction: $-0.11$, CI $[-1.36, +1.23]$ (null).
\end{itemize}

\paragraph{Additivity.} The encoder effect is $+8.37$ under CE and $+8.26$ under the full
objective --- only $\mathbf{0.11}$ pp apart. This is a particularly compelling
demonstration that the encoder contribution is additive and does not depend on the
objective.

% ---------------------------------------------------------------------
\section*{4. Isolated threshold-placement effect: $+5.68$ pp}

Tree split thresholds versus matched equal-frequency (quantile) thresholds, at
\textbf{identical channel budget and identical objective} --- only the edge placement
differs \textup{[verified]}:

\begin{quote}
Replacing the teacher's split thresholds with equal-frequency quantile thresholds at
matched budget costs $\mathbf{5.68}$ macro-F1 on the four primary datasets, 95\% CI
$[+2.31, +9.23]$ \textup{[reproduced: \texttt{scripts/reproduce\_crossed\_2x2.py},
\texttt{threshold\_placement}; historical $[+2.34,+9.26]$]}.
\end{quote}

\noindent This is distinct from both the $2\times2$ encoder effect and the
\texttt{raw\_full} sensitivity arm, and it is the \textbf{cleanest evidence about
\emph{where} inside the encoder the improvement originates}: the threshold placement
itself, not merely the existence of binning or the channel allocation. Channel budget
match verified (e.g.\ TON-IoT: 70 channels in both arms). Per dataset: CICIoT2023
$+10.77$, Edge-IIoTset $+6.50$, TON-IoT $+3.72$, CICIoMT2024 $+1.74$.

\noindent Scope: the $+5.68$ figure is the \textbf{four primary datasets}. Across all
eleven the mean is $+4.07$ pp \textup{[verified]}. \textbf{The old $+7.17$ / ``90\% of the
gain'' claim remains retired.}

% ---------------------------------------------------------------------
\section*{5. False-positive loss term: the paired effect, not just the descriptive FAR}

The inferential result is stronger than the descriptive one and should lead
\textup{[reproduced: \texttt{scripts/reproduce\_loss\_ablations.py}]}:

\begin{quote}
Removing $\alpha_{\mathrm{fp}}$ increases the false-alarm rate by $\mathbf{+2.30}$ pp,
95\% CI $[+0.06, +4.68]$, while macro-F1 changes by only $+0.15$ pp,
CI $[-0.71, +0.87]$.
\end{quote}

\noindent That is the evidence supporting the statement that the objective primarily
controls false alarms rather than aggregate accuracy. The descriptive companion (median
FAR $0.41\% \rightarrow 2.13\%$, a factor of $5.2$) is useful but secondary.

% ---------------------------------------------------------------------
\section*{6. TON-IoT protocol effect, exact model-by-model values}

Means $\pm$ SD across the four outer-fold means, five optimisation seeds within each fold
\textup{[verified --- all 15 values recomputed and matched exactly]}:

\begin{center}\small
\begin{tabular}{lrrr}\toprule
model & $\Delta$macro-F1 & $\Delta$MCC & $\Delta$FAR \\\midrule
GBDT           & $-9.34 \pm 4.63$ & $-17.21 \pm 8.15$  & $+19.32 \pm 13.90$ \\
Vanilla SNN    & $-8.13 \pm 4.27$ & $-14.91 \pm 7.62$  & $+11.95 \pm 7.46$ \\
Tree-to-Spike  & $-8.38 \pm 4.38$ & $-15.58 \pm 8.40$  & $+16.73 \pm 9.32$ \\
DNN-Tree       & $-9.90 \pm 6.07$ & $-18.18 \pm 11.26$ & $+19.25 \pm 13.99$ \\
TASNN          & $-7.97 \pm 3.58$ & $-14.93 \pm 7.06$  & $+12.06 \pm 9.01$ \\
\bottomrule\end{tabular}\end{center}

\noindent The exact table \emph{strengthens} the ``do not rank architectures''
conclusion: the degradation ranges overlap heavily, and the standard deviations are of the
same magnitude as the between-model differences.

% ---------------------------------------------------------------------
\section*{Robustness (sensitivity sweep, 700/700 cells)}

Integrity: 700 unique configurations, 700 valid records, 0 duplicate keys, 0 missing
keys. 35 arms $\times$ 4 datasets $\times$ 5 seeds, completed in 61.53 h with 0
failures. Comparisons are paired on (dataset, seed) against the shipped default run;
dataset is the generalisation unit and seeds are within-configuration optimisation
repetitions. The 700 cells are never treated as independent samples.

Default reference (median over 4 datasets $\times$ 5 seeds): macro-F1 $89.64$,
MCC $96.99$, FAR $0.41\%$, $8{,}236$ synaptic operations per sample, spike
density $10.36\%$.

\begin{quote}
Across 25 hyperparameter perturbations spanning timesteps, channel budget, pruning
quantile, jitter, distillation temperature and all four loss weights, no setting
changed macro-F1 by more than $2.06$ pp, $68\%$ of perturbations stayed within
$\pm1$ pp of the default, and every paired 95\% interval overlapped zero except the
encoder ablations. The default configuration therefore lies on a broad plateau rather
than at an isolated optimum, which supports stability around the chosen design rather
than post hoc hyperparameter optimisation.
\end{quote}

\noindent\textbf{Reporting constraint.} The nominally best arms --- \texttt{T80}
$+1.33$, \texttt{T40} $+1.23$, \texttt{epochs30} $+1.18$ pp --- all have intervals
including zero and \textbf{must not be presented as improvements} or as grounds for
changing the default.

\paragraph{Retain $T=20$.} \texttt{T40} and \texttt{T80} cost roughly $2.4\times$ and
$4.0\times$ the SOPs for nominal gains whose intervals include zero. That makes
$T=20$ a sensible cost--performance choice without claiming it is mathematically
optimal.

% ---------------------------------------------------------------------
\section*{The objective controls false alarms, not macro-F1}

\begin{quote}
The objective's practical contribution is \emph{not} an aggregate macro-F1 mechanism.
Removing the false-positive-aware term (\texttt{no\_fp}) changes macro-F1 by only
$+0.15$ pp while median FAR rises from $0.41\%$ to $2.13\%$, a factor of $5.2$.
\end{quote}

\noindent An ablation table reporting macro-F1 alone would make the objective look
useless. \textbf{Report macro-F1 and FAR side by side} (Fig.~\ref{fig:objfar},
Tab.~\ref{tab:sensitivity}). Class-weighted variants show the same pattern: FAR
$2.24$--$3.16\%$ at essentially unchanged macro-F1.

% ---------------------------------------------------------------------
\section*{Efficiency opportunity --- do NOT adopt as the deployed default}

\begin{quote}
Increasing sparsity regularisation to $10^{-2}$ reduced synaptic operations to
approximately $45.5\%$ of the default with no detectable macro-F1 loss
($-0.15$ pp) in the sensitivity study, suggesting a promising operating point for
future validation-selected deployment optimisation.
\end{quote}

\noindent This operating point was discovered \emph{using the completed sensitivity
sweep}. Adopting it as the final deployed model would be post hoc test-set selection.
It is an efficiency opportunity for future work, not a result of this manuscript.

% ---------------------------------------------------------------------
\section*{Scope qualifier (state this explicitly wherever these numbers appear)}

The sensitivity reference values, the \texttt{raw\_full} corroboration and the
\texttt{fp32} quantisation comparison all come from the \textbf{four-dataset
sensitivity subset} (CICIoT2023, CICIoMT2024, Edge-IIoTset, TON-IoT), \emph{not} the
full eleven-dataset study.

For quantisation, the supportable claim is exactly:
\begin{quote}
Across the four sensitivity datasets, int8-grid projection did not produce a
detectable Tree-to-Spike macro-F1 loss ($+0.17$ pp, 95\% CI $[-0.70, +1.90]$).
\end{quote}
It does \textbf{not} establish this across all eleven datasets, because floating-point
checkpoints are unavailable for the other seven: DNN-Tree and Tree-to-Spike run with
\texttt{quantize:\,true}, which projects weights before evaluation and saves the
checkpoint afterwards. Verified in weight space: DNN-Tree $55/55$ checkpoints exactly
on the int8 grid, Tree-to-Spike $54/55$ (one at $8.57\times10^{-8}$ relative, a
float32 residue), Raw MLP $0/55$.

By contrast the Raw MLP, trained without projection, loses $-0.97$ pp on average
(worst $-9.68$ pp; $12/55$ cells above $1$ pp), so int8-grid projection is
\textbf{dataset-sensitive} and is never described as essentially free.

Terminology throughout: \textbf{int8-grid weight projection with floating-point
execution}, never integer-only inference.

% ---------------------------------------------------------------------
\section*{Deployment --- frozen wording}

\begin{quote}
Lower analytical arithmetic count does not imply lower wall-clock latency in the
present software implementation. Sequential timestep execution dominates batch-one
CPU latency for Tree-to-Spike.
\end{quote}

\begin{quote}
Validation-selected early exit reduces executed spike operations but does not reduce
CPU latency in the current implementation; per-timestep confidence evaluation and
control overhead outweigh the saved simulation steps at batch size one.
\end{quote}

\noindent Measured: DNN-Tree $8{,}471$ MAC / $0.226$ ms; Tree-to-Spike full
$7{,}587$ SOP / $2.909$ ms; Tree-to-Spike selected exit $6{,}091$ SOP / $3.199$ ms.
\textbf{Tree-to-Spike must never be presented as the fastest implementation.} Label
the two rows \texttt{Tree-to-Spike, full $T$} and \texttt{Tree-to-Spike,
validation-selected exit}, and state on the latter that it lowers analytical SOPs but
\emph{increases} measured batch-one CPU latency.

The trade-off against the ensemble is stated as:
\begin{quote}
The full GBDT achieves the highest macro-F1, whereas Tree-to-Spike provides
substantially smaller model storage together with higher MCC and lower FAR, at the
cost of substantially higher software CPU latency.
\end{quote}

\noindent Deployment strengths: $53.4$ KB total artifact including encoder tables,
against $2.41$ MB (compact GBDT) and $5.04$ MB (full GBDT); best aggregate MCC
$90.51$; lowest FAR $3.15$; macro-F1 essentially tied with DNN-Tree.

\paragraph{Raw MLP control.} At matched int8-grid precision, Raw MLP $76.23$ vs
DNN-Tree $79.86$ vs Tree-to-Spike $79.80$: a $3.63$ pp gap showing the tree-derived
representation contributes beyond the compact dense topology alone. Raw MLP loses on
$9/11$ datasets to Tree-to-Spike and $10/11$ to DNN-Tree. \textbf{This is a
deployment/control baseline, not a formally isolated causal encoder effect} --- it
differs in representation \emph{and} training setup. The crossed $2\times2$ remains
the causal evidence.

\paragraph{CSE-CIC-IDS2018 heterogeneity.} Raw MLP exceeds DNN-Tree by $\approx6.1$ pp
there while Tree-to-Spike still exceeds Raw MLP. Per-seed DNN-Tree macro-F1 is
$48.57, 63.02, 76.44, 66.91, 74.72$ (spread $27.9$ pp) against Tree-to-Spike
$74.49$--$76.86$ (spread $2.4$) and Raw MLP $71.54$--$72.66$ (spread $1.1$). Report as
dataset-specific heterogeneity driven by DNN-Tree optimisation instability; invent no
mechanism.

% ---------------------------------------------------------------------
\section*{Early exit --- frozen wording}

\begin{quote}
Validation-selected early exit provides substantial conditional savings on datasets
for which a safe operating point exists, while automatically reverting to full-window
inference on nearly half of dataset--seed cells.
\end{quote}

\noindent Disabled in $26/55$ cells ($47\%$); four datasets reject it outright
(CICIoT2023, Edge-IIoTset, NSL-KDD, UNSW-NB15, all $5/5$). Where it qualifies:
$15$--$68\%$ fewer SOPs for $\leq0.45$ pp mean macro-F1. Only $1/55$ cells exceeded
the predeclared $1$ pp margin on test, so validation selection transfers almost
perfectly. Selected thresholds concentrate at $0.95$--$0.99$; \textbf{the retired
fixed $\theta=0.9$ headline must not reappear} (on NSL-KDD it would cost $\approx9$ pp
and no threshold qualifies at all).

Savings come from benign traffic: benign exits at step $8.02$ with $88.8\%$ exiting
early and $0.17\%$ error among early exits; attacks exit at $11.25$ with $70.9\%$
early and $16.58\%$ error. \textbf{Keep the minority-class caveat visible:} TON-IoT
ransomware $-3.83$ pp recall, CSE-CIC-IDS2018 dos $-2.29$, CICIoMT2024 MQTT $-1.85$.

Calibration justifies the framing --- ECE $23.32\%$ at $t{=}1$, $12.05\%$ at
$t{=}10$, $3.65\%$ at $t{=}20$ --- so the threshold is an adaptive-computation knob,
not a calibrated confidence level.

% ---------------------------------------------------------------------
\section*{Protocol validity I --- CTU-13 scenario-disjoint evaluation}

\paragraph{Protocol construction.} 801{,}132 flows across 13 scenarios, botnet
prevalence $55.5\%$. Timestamp audit passes: 801{,}132/801{,}132 parsed, 0 invalid, all
13 scenarios sort to valid chronological sequences, no scenario pair overlaps in time,
1\,\textmu s resolution. Four scenario-disjoint folds from a MIP under a preregistered
gate ($\geq$1000 of each class in test, $\leq60\%$ single-scenario share, 3--4 test
scenarios, $\geq$6 training scenarios, $\geq$500 of each class in validation); all four
pass and scenario overlap between train/val/test is \textbf{0} in every fold. Folds are
deliberately unequal because scenario 8 alone holds $26.8\%$ of all flows.

\paragraph{Three-condition diagnostic (RandomForest, 36 runs), means over 4 folds.}
\begin{center}\small
\begin{tabular}{lrrrrr}\toprule
features & legacy rand. & causal rand. & causal scen. & $\Delta_{\mathrm{agg}}$ &
$\Delta_{\mathrm{scen}}$ \\\midrule
base F1      & 99.99 & 99.99 & 99.23 & $+0.00$ & $-0.75$ \\
base FAR     &  0.02 &  0.02 &  1.12 & $+0.00$ & $+1.10$ \\
host F1      & 100.00 & 95.81 & 70.69 & $-4.19$ & $-25.12$ \\
host FAR     &  0.00 &  4.64 & 12.12 & $+4.64$ & $+7.47$ \\
combined F1  & 100.00 & 99.31 & 77.77 & $-0.69$ & $-21.54$ \\
combined FAR &  0.00 &  0.04 &  9.50 & $+0.04$ & $+9.47$ \\
\bottomrule\end{tabular}\end{center}

\noindent Approved wording, use verbatim:
\begin{quote}
Proper partition-local causal aggregation modestly reduces the random-split performance
of the host-context representation, but does not remove the near-ceiling result. Under
scenario-disjoint evaluation, however, host-context features become substantially less
transferable and can reduce performance when combined with otherwise stable base flow
features.
\end{quote}
Describe this as \textbf{consistent with host-context shortcut reliance}, not as a proven
mechanism.

\paragraph{$\!$Report per fold, never the mean.} \texttt{causal\_scenario} combined
macro-F1 by fold: fold 0 (sc.\ 2,4,8) $99.50$; fold 1 (sc.\ 7,9,11) $\mathbf{27.40}$;
fold 2 (sc.\ 1,3,5,10) $97.63$; fold 3 (sc.\ 0,6,12) $86.54$. The range is
$27.40$--$99.50$ and \textbf{the $77.77$ four-fold mean must never be quoted as a point
estimate}.

\paragraph{Fold-1 integrity and seed stability.} All integrity checks pass, so the
collapse is not a defect: manifest row and per-class counts match exactly, partitions
mutually disjoint, 0 NaN and 0 inf cells checked \emph{before} the stateless cleanup,
metadata disjoint from every feature list, no fitted preprocessing, and the rebuilt
feature blocks are bit-identical to the pre-fix artifact. Only 2 of 210{,}944 test rows
have an exact feature duplicate in train. Over 30 runs (5 seeds), \texttt{causal\_scenario}
combined is $29.93 \pm 8.13$ (range $24.55$--$44.19$) and host-only $34.72 \pm 21.24$
(range $24.61$--$72.71$); it never recovers in any seed.

\paragraph{Mechanism --- prediction inversion, not one hard scenario.}
\begin{center}\small
\begin{tabular}{lrrrrr}\toprule
test scenario & class balance & base F1 & combined F1 & host F1 & combined FAR \\\midrule
7  & 92.2\% Normal & 99.14 & $\mathbf{41.77}$ & 37.45 & $\mathbf{56.20}$ \\
9  & 87.0\% Botnet & 99.52 & $\mathbf{17.04}$ & 27.83 & 0.13 \\
11 & 22.1\% Botnet & 79.34 & 91.04 & 96.16 & 0.00 \\
\bottomrule\end{tabular}\end{center}
Host context fails in \emph{opposite directions}: Botnet for $54$--$66\%$ of a
$92\%$-Normal scenario, Normal for $\sim99.9\%$ of an $87\%$-Botnet scenario (DR $0.12$).
Seed-42 errors split as scenario 7 $=36{,}634$ \emph{all false positives}, scenario 9
$=106{,}233$ \emph{all false negatives}, scenario 11 $=856$.

\paragraph{$\!$The protocol is NOT host-disjoint.} $32$ of $42$ fold-1 test hosts
($76.2\%$) also appear in training, and $99.7$--$100\%$ of each scenario's flows come from
overlapping hosts. Each scenario has only $\sim$23--28 unique source hosts and
\textbf{every one is $\geq99\%$ label-pure}, so the host block acts as a $\sim$24-way host
classifier while the label is a host property. The same \texttt{147.32.84.x} addresses
recur across scenarios in different roles, so the learned host$\rightarrow$label mapping
inverts. Top-5 hosts carry $49.4\%$ of all errors ($147.32.84.170$ alone contributes
$25{,}583$ false positives). \textbf{Always call this scenario-disjoint CTU-13
evaluation, never host-disjoint.}

\paragraph{Host overlap does not explain the spread.} Fold 2 has \emph{complete}
($100.00\%$) host overlap yet stays near ceiling ($97.63$); fold 1 has the \emph{least}
($76.19\%$) yet collapses ($29.93$). The mechanism is host \emph{role inversion} plus
feature shift, not host recurrence. Host-block distribution shift is $11$--$14\times$ the
base block (mean train-standardised Wasserstein $1.837$ and $2.413$ on scenarios 7 and 9
vs $0.168$ for base); \texttt{causal\_src\_unique\_dports} has SMD $19.88$ with $56.8\%$
of scenario-7 values outside the training range.

\paragraph{Counter-cases --- host context is not uniformly harmful.} Base is $\geq99.5$ on
\textbf{12 of 13} scenarios. Host context hurts on 9, 7, 12, 3, 8, 5 but \emph{helps} on
scenario 10 ($81.47 \rightarrow 99.67$, base FAR $41.18 \rightarrow 0.95$ --- the only
scenario where base fails) and scenario 11 ($79.34 \rightarrow 91.04$). The defensible
claim is \textbf{unstable sign across unseen scenarios}, which is worse for deployment
than a uniform penalty.

\paragraph{Model matrix, 420/420 runs, 0 failures, 18.8 h.}
$\Delta_{\mathrm{aggregation}}$ (combined): GBDT $-0.11$, Vanilla $-0.31$, Tree-to-Spike
$-1.37$, DNN-Tree $-0.49$, TASNN $-0.29$; \texttt{legacy\_random} is $100.00$ for four of
five models, so \textbf{the random-split ceiling survives causalisation} and is not
explained by future-window aggregation.
$\Delta_{\mathrm{scenario}}$: GBDT $-8.32$, Vanilla $-7.68$, Tree-to-Spike $-19.24$,
DNN-Tree $-6.95$, TASNN $-8.74$ --- \textbf{every architecture degrades}.
\textbf{Do not rank them:} paired by fold, Tree-to-Spike minus GBDT is
$-0.22, -36.53, -12.68, \mathbf{+5.74}$; the sign flips on fold 3 with $n=4$ folds and
sd $23.7$, the same trap as the TON-IoT pilot ordering that dissolved.

\paragraph{Matched-mask host-context interaction (final CTU-13 computation).} Inference
only; reconstruction verified at \textbf{0.0000 pp drift on all 240 cells}. One common
mask per (condition, fold), $\mathcal{M}=\neg(\mathrm{dup}_{\mathrm{base}} \vee
\mathrm{dup}_{\mathrm{combined}})$, so both representations are scored on identical rows.
$\Delta_{\mathrm{host}}^{\mathrm{matched}}$: DNN-Tree $\mathbf{-7.90}$, GBDT
$\mathbf{-7.01}$, Tree-to-Spike $\mathbf{-9.70}$ --- negative for all three, so the
interaction is \textbf{robust}. Quote these, not the raw $-4.06/-7.98/-16.36$ or the
superseded own-mask bound $-2.97/-7.86/-13.85$, which was wrong in \emph{both} directions.
Matched evaluation collapses the between-architecture spread from $12.3$ pp to
$\mathbf{2.7}$ pp, further evidence against any ranking. Per-fold signs still flip
(DNN-Tree $+7.75$ on fold 3). On the common mask, dropping the host block gains
$\mathbf{7}$--$\mathbf{11}$ pp under scenario shift (base scenario $94.50$--$99.53$ vs
combined $83.68$--$92.47$) while costing $\sim$0--1 pp under random splitting.

\paragraph{$\!$Base-representation duplicate leakage --- state before quoting base.} With
only 8 coarse flow features, distinct flows collide: $15.5\%$ of base test rows are exact
repeats of training rows on average, rising to $\mathbf{45.9\%}$ on fold 1, and removing
them costs up to $\mathbf{17.10}$ pp (Tree-to-Spike fold-1 base $95.08 \rightarrow 77.98$).
Combined has $\sim$0\%. Base-only scenario-disjoint scores are therefore inflated.

\paragraph{Per-scenario, architecture $\times$ scenario specific and inverting.}
\begin{center}\small
\begin{tabular}{lrrrrrrr}\toprule
scenario & botnet\% & GBDT & Vanilla & Tree-to-Spike & DNN-Tree & TASNN & spread \\\midrule
7  &  7.8 & $\mathbf{98.20}$ & 39.35 & 56.50 & 41.24 & 38.60 & $\mathbf{59.6}$ \\
9  & 87.0 & 61.05 & $\mathbf{99.73}$ & $\mathbf{27.18}$ & 99.38 & 98.87 & $\mathbf{72.6}$ \\
10 & 75.0 & 99.42 & 96.20 & 60.85 & 65.39 & 94.71 & 38.6 \\
12 & 55.6 & 73.31 & $\mathbf{99.92}$ & 84.06 & 96.24 & 87.74 & 26.6 \\
1  & 69.7 & 100.00 & 78.84 & 78.64 & 96.95 & 85.47 & 21.4 \\
\bottomrule\end{tabular}\end{center}
The other 8 scenarios are $\geq93$ for every model. \textbf{No architecture is uniformly
robust}: GBDT is near-perfect on scenario 7 (FAR $0.31$) where all four neural models
collapse (FAR $38.7$--$60.7$), and is among the worst on 9 and 12. The strongest CTU-13
claim is therefore:
\begin{quote}
Near-perfect random-split performance conceals scenario-specific prediction inversions,
and no architecture is uniformly robust across held-out scenarios.
\end{quote}

% ---------------------------------------------------------------------
\section*{Protocol validity II --- TON-IoT source-endpoint-disjoint evaluation}

\begin{quote}
Under source-IP-disjoint evaluation all five architectures lose $8$--$10$ macro-F1 and
gain $12$--$19$ pp FAR. The effect is \textbf{protocol-wide degradation, not an
architecture ranking}.
\end{quote}
A pilot with one seed per cell showed Tree-to-Spike degrading least; with 5 seeds per fold
that ordering \textbf{dissolved}. \textbf{Do not revive the ranking claim.} Package frozen
at 200/200 runs, 0 failures, with 8 frozen fold manifests.

\paragraph{Protocol feasibility.} Only \textbf{1 of 4} target-domain datasets supports a
new group-disjoint protocol (TON-IoT). CICIoMT2024 has provider capture separation;
CICIoT2023 and Edge-IIoTset are infeasible (documented in \texttt{rebuttal2\_response.tex},
\texttt{tab:audit}).

% ---------------------------------------------------------------------
\section*{CICIoMT2024 exact-record overlap (leakage audit subsection, not appendix)}

Required wording:
\begin{quote}
The provider-defined CICIoMT2024 partitions are capture-file-disjoint at the filename
level, but exact-record matching identifies 86{,}012 test records also represented in
training, corresponding to $5.33\%$ of the full provider test set. Under the study's
seed-specific $25\%$ sampling protocol, approximately $5.2\%$ of each evaluated test
subset overlaps its corresponding sampled training subset. All overlapping records have
consistent labels. Removing them changes macro-F1 by at most $0.36$ pp for any model and
leaves the benign false-alarm rate exactly unchanged, but the tree-derived models score
$\sim99.9\%$ on the overlapping records against $\sim89$--$90\%$ on unique ones, so the
overlap is being memorised even though its effect on the aggregate metric is small.
\end{quote}

\noindent Per-seed removal $20{,}973$--$21{,}204$ of $403{,}545$ ($5.224\%$ mean); label
agreement $86{,}012/86{,}012 = 100\%$. Duplicate-free re-evaluation was inference only
from existing checkpoints, provider partition unchanged, 25 evaluations.

\begin{center}\small
\begin{tabular}{lrrrrr}\toprule
model & macro-F1 before $\rightarrow$ after & $\Delta$F1 & $\Delta$MCC & dup.\ rows &
unique rows \\\midrule
GBDT          & $90.541 \rightarrow 90.540$ & $-0.001$ & $-0.010$ & $\mathbf{99.896}$ & 90.540 \\
Tree-to-Spike & $89.424 \rightarrow 89.419$ & $-0.005$ & $-0.013$ & $\mathbf{99.867}$ & 89.419 \\
DNN-Tree      & $89.240 \rightarrow 89.234$ & $-0.006$ & $-0.016$ & $\mathbf{99.882}$ & 89.234 \\
Vanilla SNN   & $72.608 \rightarrow 72.874$ & $+0.266$ & $\mathbf{+1.204}$ & 61.293 & 72.874 \\
TASNN         & $74.240 \rightarrow 74.581$ & $+0.341$ & $\mathbf{+1.404}$ & 61.428 & 74.581 \\
\bottomrule\end{tabular}\end{center}

\noindent$\Delta$FAR is \textbf{exactly $0.000$ for every model} because zero Benign rows
are duplicated. Model ranking is unchanged. Removals concentrate in DDoS ($15{,}167$) and
DoS ($5{,}848$) with Recon $68$; \textbf{Benign, MQTT and Spoofing lose exactly zero rows},
so the weakest class (Spoofing, F1 $40.4$) is untouched.

\paragraph{Verdict.} Retain the provider capture-file split as the primary predictive
result and promote this duplicate-free analysis to the main leakage-audit subsection. Do
not make duplicate-free the preferred result: no model moves $>1$ pp macro-F1 and no
minority-class F1 shifts.

\paragraph{TASNN caveat.} Four of five models reconstruct bit-exactly (0.0000 pp drift);
\textbf{TASNN does not} (up to $+0.2386$ pp in 3 of 25 cells), which is close to its own
measured $\Delta$ of $+0.341$. \textbf{Report TASNN's structural overlap but do not quote
its exact duplicate-free delta.} Vanilla SNN establishes the same opposite-sign effect at
zero drift.

% ---------------------------------------------------------------------
\section*{Compact GBDT --- a validation-selected baseline, not a guarantee}

Validation-only tree-prefix selection (doubling grid then binary search for the true
smallest qualifying prefix); all 55 cells found one. Median round compression
$\mathbf{2.74\times}$ (the $25.9\times$ mean is an artifact of CTU-13 needing a single
round), $53.7\%$ fewer nodes overall.

\noindent\textbf{Disclose the overshoot:} $\mathbf{20/55}$ cells lose more than $1.0$ pp
on \emph{test} (4 lose $>2$ pp; worst UNSW-NB15 $-2.01$) despite every cell being selected
to lose at most $1.0$ pp on validation (max validation drop $-0.998$ by construction).
Describe it as a \textbf{validation-selected compressed baseline}, never as a model
guaranteed to stay within 1 pp on test.

\paragraph{Informative contrast.} Under the same nominal tolerance, validation-selected
temporal early exit exceeded the margin in only $\mathbf{1/55}$ cells against compact
GBDT's $\mathbf{20/55}$. Report the empirical difference; do not overstate why.

% ---------------------------------------------------------------------
\section*{Reproducibility --- environment of record}

\begin{center}\small
\begin{tabular}{ll}\toprule
Python 3.12.3 & torch 2.12.0+cu130 (CUDA 13.0) \\
numpy 2.4.4 & pandas 3.0.2 \\
scikit-learn 1.8.0 & scipy 1.17.1 \\
xgboost 3.2.0 & snntorch 0.9.4 \\
\bottomrule\end{tabular}\end{center}

\noindent Hardware: 1$\times$ NVIDIA H200 NVL (143 GB), 16 CPU cores, 125 GB RAM.
Sensitivity sweep completed in $61.53$ h with 0 failures; CTU-13 matrix in $18.8$ h with
0 failures; frozen fold manifests are hashed (CTU-13 scenario manifest
\texttt{sha256=79ec75d7\ldots}) and the hash is recorded into every result file produced
against it.

% ---------------------------------------------------------------------
\section*{Multiplicity: the two nominally significant loss ablations do NOT survive}

Isolated ablation contrasts used \emph{unadjusted} 95\% intervals. Two are nominally
significant \textup{[reproduced: \texttt{scripts/reproduce\_loss\_ablations.py}]}:
ranking-only $[-2.39, -0.29]$ ($p_{\mathrm{raw}}=0.0044$) and class-weighted-no-KD
$[-2.46, -0.05]$ ($p_{\mathrm{raw}}=0.0418$). After correction over the declared family
of $m=16$: $p_{\mathrm{Bonf}}=0.070$ and $0.669$; $p_{\mathrm{BH}}=0.070$ and
$0.334$. \textbf{Neither survives Bonferroni or Benjamini--Hochberg correction, and
nothing in the family does.}

\noindent Family membership is \emph{declared in the script}, not inferred: every arm
perturbing only loss weights or class weighting, which gives $m=16$. Both nominally
significant contrasts must sit in one family or the multiplicity claim is untestable.
$m=16>11$ makes both corrections \emph{stricter}, so the conclusion is conservative ---
nothing failing at $m=16$ would pass at $m=11$.

\begin{quote}
Every isolated ablation macro-F1 effect is at most $1.28$ pp in
magnitude and small relative to the encoder effect.
\end{quote}

\noindent This is recorded explicitly so a future manuscript edit cannot accidentally
resurrect either nominally significant result as a finding.

% ---------------------------------------------------------------------
\section*{Class-weighting control (direct reviewer concern)}

\begin{itemize}
  \item Inverse-frequency class weighting \emph{on top of} distillation: macro-F1
        $\mathbf{-0.16}$ pp, CI $[-1.43, +1.27]$, at $109.3\%$ of default spiking
        operations \textup{[reproduced]}. \textbf{The historical ``about $+5\%$''
        operation increase does not reproduce as a single figure}: measured per arm it is
        $+9.3\%$ (\texttt{class\_weighted}), $+10.5\%$ (\texttt{class\_weighted\_nokd})
        and $+5.6\%$ (\texttt{cw\_nokd\_rankfixed}). Quote the arm, not an average ---
        only \texttt{cw\_nokd\_rankfixed} is near $+5\%$.
  \item Replacing distillation with class-weighted CE: $\mathbf{-1.28}$ pp,
        CI $[-2.46, -0.05]$ \textup{[reproduced]} --- but see the
        multiplicity section, this does not survive correction.
\end{itemize}

\begin{quote}
The encoder and distillation behaviour cannot be explained as indirect transfer of the
teacher's class weighting.
\end{quote}

% ---------------------------------------------------------------------
\section*{Temporal-information capacity at $T=20$ (reviewer-requested, now all 11 datasets)}

\texttt{scripts/encoder\_temporal\_capacity.py} $\rightarrow$
\texttt{encoder\_temporal\_capacity.json}. Inference only. The computation
\textbf{reproduces the committed four-dataset artifact to floating-point precision}
(worst absolute difference $1.11\times10^{-16}$) and is then extended to all eleven
datasets \textup{[verified]}.

A feature $j$ with $|E_j|$ split thresholds has $|E_j|+1$ reachable bin levels, while the
latency code offers at most $T$ distinct spike times, so whenever $|E_j|+1 > T$ distinct
split outcomes must share a spike time.

\begin{center}\scriptsize
\begin{tabular}{lrrrrrr}\toprule
dataset & kept feat. & med.\ $|E_j|$ & max $|E_j|$ & $T\!\geq\!|E_j|{+}1$ & collision &
recoverable \\\midrule
CTU-13          &  22 &  15.5 &  34 & 86.4\% & \textbf{7.6\%}  & \textbf{92.4\%} \\
KDD Cup 99      &  56 &   2.0 & 128 & 69.6\% & 59.9\% & 40.1\% \\
CIC-IDS2017     &  54 &  38.0 & 152 & 29.6\% & 65.2\% & 34.8\% \\
NSL-KDD         &  61 &   1.0 & 167 & 68.9\% & 65.3\% & 34.7\% \\
ISCX-IDS2012    &  49 &   1.0 & 170 & 87.8\% & 69.7\% & 30.3\% \\
Edge-IIoTset    &  37 &   1.0 & 234 & 83.8\% & 77.0\% & 23.0\% \\
CSE-CIC-IDS2018 &  57 &  70.0 & 187 & 24.6\% & 77.4\% & 22.6\% \\
TON-IoT         &  55 &   1.0 & 238 & 81.8\% & 78.5\% & 21.5\% \\
UNSW-NB15       & 152 &   1.0 & 254 & 82.2\% & 78.8\% & 21.2\% \\
CICIoMT2024     &  33 &  87.0 & 237 &  9.1\% & 81.3\% & 18.7\% \\
CICIoT2023      &  31 & 103.0 & 255 & 29.0\% & \textbf{87.1\%} & \textbf{12.9\%} \\
\bottomrule\end{tabular}\end{center}

\paragraph{$\!$Scope correction.} The previously reported ``$77$--$87\%$ of reachable
levels collided at $T=20$'' is correct \emph{for the four primary datasets only}
($77.0$--$87.1\%$). Across all eleven the range is much wider,
$\mathbf{7.6}$--$\mathbf{87.1\%}$ (median $77.0\%$), and CTU-13 is a strong outlier at
$7.6\%$ because its teacher uses at most 34 thresholds per feature. \textbf{Do not
generalise the $77$--$87\%$ figure to all eleven datasets.}

\noindent The substantive point survives and is reportable: on the primary target-domain
datasets the TTFS code at $T=20$ discards most of the teacher's threshold resolution
($12.9$--$23.0\%$ recoverable), yet Tree-to-Spike still realises the encoder gain --- so
the benefit does not require faithful reconstruction of every split outcome. This is
consistent with the sensitivity result that $T=40$ and $T=80$ buy nothing detectable.

% ---------------------------------------------------------------------
\section*{Two smaller items worth freezing}

\paragraph{Spike sparsity, primary set: $0.767$--$0.919$ \textup{[verified]}.} State the
basis explicitly: this is the \textbf{range of per-dataset means} over the four primary
datasets (TON-IoT $0.767$, CICIoT2023 $0.862$, Edge-IIoTset $0.911$, CICIoMT2024
$0.919$). The per-\emph{run} range on the same set is wider, $0.746$--$0.926$, and the
range of dataset means across all eleven is $0.767$--$0.972$. Quoting one basis as another
is the likely source of any future mismatch.

\paragraph{TASNN scope.} The evaluated shared-input TASNN adaptation is substantially
larger than the compact students and is potentially under-optimised. \textbf{All TASNN
conclusions apply only to that adaptation, not to transformer-SNNs in general.} This
qualifier must accompany the $+6.03$ pp figure in Section 2 and every TASNN row in the
deployment and protocol tables.

% ---------------------------------------------------------------------
\section*{$\!$Caution: regenerate the analytical-arithmetic figures}

\textbf{Do not copy the pre-deployment $10.5$K pJ / $39.2$K pJ / $3.7\times$ arithmetic
figures into this file or the manuscript.} The measured operation counts changed
substantially during the deployment audit, so the definitive analytical-arithmetic table
must be generated from the final \texttt{deployment\_summary.csv} operation accounting
(7{,}587 SOP full / 6{,}091 SOP selected exit / 8{,}471 MAC DNN-Tree /
10{,}979 MAC Raw MLP), with tree models priced in comparisons and never converted to
MAC-equivalents. Label the result \emph{analytical arithmetic cost}, not energy, and keep
32-bit pricing as the principal estimate with int8 repricing as sensitivity only.

% ---------------------------------------------------------------------
\section*{Provenance --- every statistic regenerates from committed results}

\texttt{results/final\_provenance.json} indexes six statistical entries, each naming the
script and inputs that regenerate it. No training was rerun; all five formerly
prose-only analyses are recomputed from committed result JSON/CSV.

\begin{center}\small
\begin{tabular}{ll}\toprule
analysis & producer \\\midrule
canonical bootstrap, BH, Bonferroni & \texttt{scripts/statistics\_utils.py} \\
dense-twin non-inferiority & \texttt{scripts/reproduce\_noninferiority.py} \\
crossed $2\times2$ attribution & \texttt{scripts/reproduce\_crossed\_2x2.py} \\
FAR term, multiplicity, class weighting & \texttt{scripts/reproduce\_loss\_ablations.py} \\
temporal capacity at $T$ & \texttt{scripts/encoder\_temporal\_capacity.py} \\
manifest assembly & \texttt{scripts/make\_provenance\_manifest.py} \\
\bottomrule\end{tabular}\end{center}

\paragraph{Bootstrap design, frozen.} Resampling unit is the \textbf{dataset}; seeds are
resampled \emph{within} each drawn dataset (nested). The statistic is the mean over
datasets of the per-dataset mean paired difference, so datasets are equally weighted
regardless of seed count. \textbf{10{,}000} replicates, RNG seed \textbf{20260812},
\textbf{percentile} intervals (not BCa --- the acceleration term is unstable at $n=4$
datasets), $\alpha=0.05$, all effects in \textbf{percentage points}, signed as
\textbf{arm minus reference}. This metadata is emitted into every result payload, because
a rerun under a different convention would produce a different-but-plausible interval
that nobody would catch.

\paragraph{Reproduction fidelity.} Every reproduced value matches its historical
counterpart to bootstrap noise: the $2\times2$ cell means match \emph{exactly}
($77.86/77.59/86.23/85.85$); NI CI $[-1.93,+1.03]$ vs $[-1.92,+1.03]$; encoder
$[+3.57,+13.08]$ vs $[+3.56,+13.00]$; MCC encoder $[+2.18,+33.80]$ vs $[+2.17,+33.79]$;
$\Delta$FAR $[+0.06,+4.68]$ vs $[+0.05,+4.72]$. \textbf{Where they differ, the reproduced
value is authoritative.} The temporal-capacity script additionally reproduces its
committed four-dataset artifact to $1.11\times10^{-16}$.

% ---------------------------------------------------------------------
\section*{Retired claims --- must not reappear anywhere}

\begin{itemize}
  \item ``Causal aggregation has no effect'' / ``CTU-13's ceiling is purely host
        identity leakage'' / ``host-only performance remains exactly at ceiling after
        causalisation.''
  \item ``CICIoMT2024 contains only a negligible $0.02\%$ duplicate overlap''
        (195 clusters / 357 rows). Measured: $86{,}012$ of $1{,}614{,}182$ full
        provider test records ($5.33\%$), $\approx5.2\%$ per evaluated seed subset,
        $100\%$ label agreement.
  \item Describing any CTU-13 protocol as \emph{host-disjoint} --- $76.2\%$ of fold-1
        test hosts also appear in training. Use \emph{scenario-disjoint CTU-13
        evaluation}.
  \item ``Host identity is a formal label leak.'' Use \emph{host-associated shortcut
        reliance}.
  \item The fixed $\theta=0.9$ early-exit headline.
  \item ``int8-grid quantisation has little effect'' as a general statement.
  \item Any unqualified four-fold mean for CTU-13 scenario-disjoint results (the
        $77.77$ mean spans $27.40$--$99.50$).
  \item Any architecture ranking from CTU-13 or TON-IoT scenario/host-disjoint
        results.
\end{itemize}

% ---------------------------------------------------------------------
\section*{Artifacts in this directory}

\begin{tabular}{ll}
\texttt{fig\_sensitivity\_forest.png} & paired $\Delta$ macro-F1 with 95\% intervals,
35 arms \\
\texttt{fig\_sensitivity\_cost\_quality.png} & SOP\% vs $\Delta$ macro-F1, default
marked \\
\texttt{fig\_objective\_far.png} & median FAR by objective ablation \\
\texttt{fig\_deployment\_pareto.png} & macro-F1 vs measured batch-one CPU latency \\
\texttt{tab\_sensitivity.tex} & 35-arm sensitivity table \\
\texttt{tab\_deployment.tex} & seven-row deployment comparison \\
\texttt{tab\_early\_exit.tex} & validation-selected early exit per dataset \\
\texttt{tab\_raw\_mlp.tex} & Raw MLP control, FP and int8-grid \\
\texttt{sensitivity\_by\_arm.csv}, \texttt{sensitivity\_cells.csv} & sweep source data \\
\texttt{deployment\_summary.csv}, \texttt{int8\_grid.csv} & deployment source data \\
\texttt{compact\_gbdt.csv}, \texttt{selected\_test.csv} & compact GBDT, early exit \\
\end{tabular}

\bigskip
\noindent Regenerate everything with
\texttt{python scripts/make\_journal\_figures.py}.
